\chapter{Who Decides Which Works Are Classics?}
\section{A Reading List at a Bookshop Door}
Pick up something you have certainly seen: a reading list.

In the first school week, your form teacher sends the parents' group a document titled "Recommended Extracurricular Reading for Middle-School Pupils". The bookshop's study-guides section promptly prints it and posts it prominently by the entrance, circling several lines in red. Look closer. Under "Classical Chinese Masterpieces": Journey to the West, Water Margin, Dream of the Red Chamber. Under "Modern and Contemporary Literature": Dawn Blossoms Plucked at Dusk, Rickshaw Boy, Border Town. Foreign classics follow: Twenty Thousand Leagues Under the Seas, How the Steel Was Tempered, Jane Eyre.

Thin paper, astonishingly dense information. Notice several things.

First, it is a list, meaning some names appear and many more do not. Journey to the West appears; Investiture of the Gods may not. Lu Xun appears; Jin Yong probably does not. Jane Eyre appears; Harry Potter certainly does not. Why?

Second, none of those names climbed there unaided. Dream of the Red Chamber was written over two centuries ago. During its author's lifetime it circulated among friends as a manuscript, its authorship itself unclear. Dawn Blossoms Plucked at Dusk consists of recollective essays written less than a century ago by a man from Shaoxing who died not knowing he would enter every Chinese pupil's textbook. Between an ordinary, perhaps nearly lost book and this printed "must-read" lies a long chain of people's decisions.

Third, this paper has power. Listed books sell hundreds of thousands more copies, enter successive generations' reading reports, and become standard answers to literary-knowledge questions. Unlisted books gradually leave bookshop shelves and library recommendation stands, finally fading from memory. A reading list divides a work's destiny.

This sheet is therefore the machine we will take apart. Our question is simple and sharp: \textbf{who decides which works are classics?} Are the works so good that history cannot forget them? Or do editors, critics, teachers, officials, and booksellers decide on history's behalf? If the latter, by what right? Have they misjudged? What happened to the books and people they kept outside?

\section[A Bell and an Ancient Call for Songs]{A Bell Passes the Village: An Ancient Call for Submissions}
Come into a scene so early that bookshops, publishers, and even the profession of author do not yet exist.

Time: some twenty-seven or twenty-eight centuries ago, a Zhou spring. Place: a dirt road outside a Central Plain village.

Winter's slack season is ending, but farm work has not grown busy. Elders gather in the sunshine at the village entrance, humming songs. Nobody "wrote" these songs; countless mouths sang them into being: young men's love songs, wives' laments about endless war, thanksgiving songs after harvest. The wind carries them far.

A man comes along the road, looking less like an official than an itinerant pedlar. He shakes a large bell with a wooden tongue, a muduo. Its ringing tells the villagers that the song collector has arrived.

The "Treatise on Food and Money" in the Book of Han, written in the Eastern Han, records a later description of this institution:

\begin{quote}
"In the first spring month, before the assembled people disperse, an envoy shakes a wooden-tongued bell along the roads to collect poems."
\end{quote}
The sense: each early spring, before people scatter into the fields, an envoy called a xingren travels the roads ringing his bell and collecting folk songs. These pass upward through successive levels; the Grand Music Master sets them to music, and they are finally sung to the Son of Heaven, informing him how people live and what they resent.

Stay at that village entrance a moment longer. It holds this chapter's central picture: \textbf{a song travelling from a field ridge into a palace.}

Some songs later formed the Book of Songs. In the Records of the Grand Historian's "Hereditary House of Confucius", Sima Qian records a claim that over three thousand ancient poems existed before Confucius removed repetitions and those incompatible with ritual and righteousness, retaining 305. This is the famous account of Confucius editing the poems. Most modern scholars doubt it, since many of the "Three Hundred Poems" circulated among aristocrats before his birth. Yet the legend reveals something: ancient people clearly understood that creating a classic requires \textbf{selection and compilation}. They simply assigned that task to their most prestigious figure.

The collection's status kept rising. Han authorities appointed specialist scholars to teach it, transforming poems into a canonical scripture. For over a millennium afterward it supplied compulsory examination material, and children across China began literacy with "Guan-guan cry the ospreys on an island in the river". The young men, wives, and veterans who originally sang at village entrances left not even their names.

How many hands intervene between anonymous field song and revered scripture? Collectors decide which songs to gather or miss; musicians decide how refined singing should sound; compilers decide what enters and where; commentators decide what it ought to mean; examiners decide who must memorise it. The songs remain songs, but this chain of people raises their status stage by stage.

That is the production line we will trace. Remember the bell's ringing.

\section{Four Lists That Are Never the Same}
Before proceeding, map the territory without rushing to conclusions.

What people call "good books" actually belongs to four different lists, often overlapping but never identical.

\textbf{List 1: Bestsellers.} One criterion: sales. Prominent shop displays and online sales rankings give this list physical form. Every paying reader judges; awards arrive weekly or even daily.

\textbf{List 2: Literary prizes.} The Nobel Prize in Literature, Mao Dun Literature Prize, American Pulitzer, British Booker. A few insiders judge annually. Besides money they award recognition, a stamp declaring serious literature.

\textbf{List 3: Syllabuses.} Textbook contents, recommended reading, examination outlines. Textbook editors and educational authorities judge on a scale of decades. This list wields the greatest power because it prescribes, not suggests: children nationwide read the same text in the same classroom.

\textbf{List 4: Classics.} The strangest list: no judges, ceremony, or clearly defined boundaries. Its loose criterion is that decades or centuries later, people still read, quote, remake, and argue about the work.

These lists diverge more than you might imagine. Today's biggest seller often vanishes from conversation within five years. Prizes have chosen wrongly: the inaugural 1901 Nobel went to French poet Sully Prudhomme, while widely favoured Tolstoy received nothing. Outraged, dozens of Swedish writers and artists jointly wrote to apologise to Tolstoy. Do you remember anything Prudhomme wrote? Syllabuses, meanwhile, lag behind: a book usually must prove sufficiently safe elsewhere before entering a classroom.

Where is the conflict? Between two opposite claims, both sounding reasonable.

Claim one: good works survive by themselves. Time is the fairest sieve, panning away sand and retaining gold. A book surviving three centuries must have reasons. If your grandfather, father, and you all read it, nobody needs to decide: it is a classic on its own merits.

Claim two: survivors are those permitted to survive. Time is no fair sieve. A book burned in war gets no second chance; a form printing disfavour leaves no chance of survival; an author despised by the mainstream may never enter the sieve. Rather than time's selection, our canon is the result of generations of gatekeepers granting passage. Gold certainly sinks, but what if it never entered the river?

These are the opponents we will repeatedly compare. For now, ask yourself: does the village bell \textbf{discover} good songs or \textbf{create} them?

\section{Gatekeepers and Those Outside}
First inspect both sides of the gate.

One side holds a line of gatekeepers. Who are they?

Nearest is the \textbf{publisher}: the Zhou collector and Music Master, Song printing workshop, today's publishing-house editor. Any book seeking many readers must first pass someone's decision that it deserves printing. Then come \textbf{critics}: newspaper book-page editors, prize juries, university professors, telling the public which books deserve serious attention. \textbf{Teachers and textbook editors} decide what the next generation must read. The \textbf{market} matters too: booksellers choose stock and placement, readers whether to buy. Finally, often forgotten yet strongest of all, comes \textbf{political power}.

How does power keep gates? Two Chinese scenes show it. First, elevation: Emperor Wu of Han established Erudites of the Five Classics, supporting scholars of the Songs, Documents, Rites, Changes, and Spring and Autumn with state money and office. Later civil-service examinations locked scholarly careers to Confucian classics, and from the Yuan even standard answers were fixed in Zhu Xi's commentaries. This was history's most powerful syllabus: controlling interpretation controlled scholars' futures. Second, suppression: compiling the Complete Library of the Four Treasuries under Qing emperor Qianlong, the state simultaneously collected and copied old books and banned, destroyed, or altered thousands for objectionable content. One state machine compiled lists with one hand and burned books with the other. Ming and Qing rulers also repeatedly banned and destroyed novels such as Water Margin for "teaching banditry": fear they would teach rebellion.

Who stands outside? Visit several gates worldwide.

In early eleventh-century Japan, Heian court lady Murasaki Shikibu wrote the long narrative The Tale of Genji. Formal writing then used literary Chinese, the male language of officialdom. Women were excluded from Chinese-language education and restricted to kana, which recorded vernacular speech. Yet precisely through this unofficial women's script, Murasaki portrayed inner life with unprecedented subtlety. Today Genji is called Japanese literature's supreme classic. A classic emerged through the gate the gatekeepers despised.

Nineteenth-century Arab scholars disdained a centuries-old popular story collection: sailors, demons, flying carpets, tales stretching night after night. Its language was colloquial, its versions untidy, and anyone might add another story. This was the Thousand and One Nights. Early in the eighteenth century, French translator Galland translated it, causing a European sensation before further translations spread worldwide. What Arab serious-literature gatekeepers rejected became children's bedside reading across the globe.

In West Africa, history on Mali's grasslands and in its villages lives not in books but in people's heads. Griots, hereditary oral historians and singers, perform the epic Sundiata for days and nights, recounting a thirteenth-century founding ruler's life. If classics must be printed, West Africa seems to have none; the measuring stick itself is wrong. Many literary historians only learned to treat oral traditions as literature in the twentieth century.

Now an important exercise: \textbf{state the gatekeepers' full case}. "Gatekeeper" carries a negative tint, but criticism becomes mere emotion unless their reasoning first receives its strongest form.

They have at least four reasons. First, \textbf{selection is necessary}. Millions of new books appear globally each year; a lifetime allows only thousands. Without selection everyone searches a rubbish dump for diamonds by hand. Second, \textbf{expert judgement has real value}. An editor who has read ten thousand books genuinely judges differently from a child just learning literacy. Admitting this is honest, not snobbish. Third, \textbf{teaching time is limited}. Six secondary-school years allow perhaps one or two hundred closely studied texts. Every choice excludes ten thousand alternatives. Someone must make that painful choice for children; repeatedly tested consensus offers the least risk. Fourth, easily overlooked, \textbf{a world without gatekeepers is not necessarily fair}. Remove editors and critics and selection power remains, transferring to the loudest, richest, most attention-savvy actors. Recommendation algorithms and trending lists are the gatekeepers of an allegedly gatekeeper-free age, yet never answer for their criteria.

Strong reasoning. But those outside hold something the gatekeepers cannot possess: \textbf{the list the sieve let fall through}. Its holes have a direction. All-male critics call women's writing delicate but narrow; juries reading only European languages collectively deem other literatures immature; a class defining refinement permanently labels others' tastes popular. The problem is not a sieve's existence, but its pretence of having no holes.

\section[Thinking Bridge: The Canon-Making Line]{A Thinking Bridge: Dismantle the Canon-Making Production Line}
Talent and intuition cannot answer why a book is a classic; we need a portable tool. Simply arrange what we have seen into a production line: \textbf{a work becoming a classic must pass five successive checkpoints, each with different gatekeepers.}

\textbf{Checkpoint 1: Be written and survive.} A work must exist and withstand war, prohibition, insects, even its author's abandonment. Luck guards this gate, with material conditions helping: enough copies, repeated copying. Many works perish here silently.

\textbf{Checkpoint 2: Be published.} Manuscript becomes printed edition, workshop becomes publishing house. Booksellers and editors guard the gate, asking whether spending money is worthwhile.

\textbf{Checkpoint 3: Be evaluated.} Critics write, prizes stamp approval, anthologists select. Intellectuals guard the gate with their tastes and judgements.

\textbf{Checkpoint 4: Be taught.} Enter textbooks, lists, examinations. Education guards the gate, asking whether a work suits the next generation. That criterion includes morality and politics as well as aesthetics.

\textbf{Checkpoint 5: Be reread across generations.} This alone has no gatekeeper: ordinary readers across time are the judges. Recommendations can secure a first reading; a tenth reading, and grandchildren continuing to read, depend on the book itself.

Use the tool in three questions: \textbf{take any recognised classic through the line and ask at each station: who guards this gate, by what standard, and who remains outside?}

Try Dream of the Red Chamber. First checkpoint: under Qianlong it circulated among friends and relatives as The Story of the Stone, surviving through fragile word-by-word manuscript copying without formal publication. Second: nearly thirty years after its author's death, booksellers Cheng Weiyuan and Gao E organised and completed it, issuing a movable-type edition in 1791. Only then did it first become a book one could buy. Third: twentieth-century scholars including Hu Shi investigated its author's family, elevating this leisure reading into academia and establishing Redology. Fourth: it entered secondary-school textbooks and examinations; the very category "four great classical novels" is a joint product of twentieth-century publishing and teaching. Fifth: it continues to be reread, remade, and debated. Nothing here happened naturally: specific people made specific decisions at every step.

The tool extends beyond books. Analyse a century-old restaurant universally deemed delicious, a supposedly best school, or a universally classic internet meme. Wherever consensus appears, a production line and usually unnamed gatekeepers stand behind it.

\section[Evidence: A Misjudged Ranking]{A Little Evidence: A Ranking That Put People in the Wrong Order}
Now some primary material: a documented literary misjudgement.

Around fifteen centuries ago, during the Southern Dynasties' Liang, critic Zhong Rong wrote Grades of Poets, ranking writers of five-character verse since the Han into upper, middle, and lower grades. Among China's earliest poetic rankings, it placed today's highly revered Tao Yuanming in the middle.

Interestingly, Zhong did recognise quality, giving Tao exceptionally high praise:

\begin{quote}
"The master of poets of reclusion, ancient and modern."
\end{quote}
Meaning: among all writers of reclusive poetry, Tao is the founding master. Yet this acknowledged master receives only middle rank because Zhong's generation prizes ornate diction and elaborate allusion. Tao's simplicity, "I plant beans beneath the southern hill" and "With the moon above, I shoulder my hoe home", fails its standard of sophistication. In other words, \textbf{Tao is not inadequate; the ruler cannot measure him}.

Look at the ruler's other end. Early Tang poets Wang Bo, Yang Jiong, Lu Zhaolin, and Luo Binwang, collectively Wang--Yang--Lu--Luo, were widely mocked for a frivolous style. Decades later Du Fu defended them in "Six Quatrains Written in Jest", including:

\begin{quote}
"Your bodies and names will perish together; the rivers will flow through the ages undiminished."
\end{quote}
The sense: you mockers will lose your names when you die, while the works you mock flow forever like rivers. It sounds prophetic, but also provides evidence: even Du Fu saw that \textbf{contemporary judgements are often unreliable}.

His own experience is especially telling. His poetic reputation was modest while alive; several contemporary Tang anthologies barely selected him. Decades after his death, eminent writer Yuan Zhen's epitaph elevated him unprecedentedly, effectively declaring that no poet since poetry began matched Du Zimei. Two centuries later Song readers made him the "Poet Sage". Between anthology exclusion and sainthood, Du Fu wrote not one further word: \textbf{the rankings changed, not the poems}.

These three pieces of evidence indicate an unsettling fact: literary rankings are not works voting for themselves but generations of readers voting with rulers in hand. Rulers change, so rankings change. Why, then, should today's ranking be right?

\section[Three Accounts of Tao's Promotion]{Comparing Explanations: How Did Tao Yuanming Move Up a Class?}
Separate four things, a discipline this book repeatedly insists upon. \textbf{What happened}: Zhong Rong ranked Tao in the middle; after the Song, Tao securely occupied the first rank. These are documented, undisputed facts. \textbf{Why}: interpretation, where accounts conflict. \textbf{How participants understood it}: Zhong sincerely regarded middle rank as praise; this is his own account. \textbf{Our evaluation}: calling Zhong mistaken is our judgement. We compare the second layer: why Tao's standing rose.

\textbf{Explanation 1: True value eventually emerges: time sifts.}

Contemporaries stand too close, blinded by passing fashion. With distance, noise settles and gold appears. Tao's depth beneath simplicity required more mature readers. Comforting and partly sound, this account recognises that works sustaining repeated reading across eras usually have rich layers. Its fatal limit: \textbf{it counts only survivors}. How many equally good or better writers disappeared through war, prohibition, or failure to be copied? They never reached emergence day. Time-sifting theory sees surfaced gold but not what remains on the riverbed.

\textbf{Explanation 2: Each age rewrites rankings to meet its needs: historical taste.}

Tao's fastest promotion came in the Northern Song. Its literary leader Su Shi was at a personal low, demoted to Huangzhou and thwarted in official life, urgently needing the posture of someone who "would not bow for five pecks of rice". He immersed himself in Tao and composed over a hundred poems matching Tao's rhymes, effectively casting his entire late-life reputation as a vote. Su did not objectively discover Tao: he \textbf{needed} a Tao. Each age's rearrangement of classics resembles looking into a mirror, finding what it lacks. A sharp account, but if literary history merely changes idols with the mood, why do some promoted authors never become beloved, while others such as Tao, after centuries of neglect, become impossible to put down once seriously read? The work itself clearly casts a vote too.

\textbf{Explanation 3: Institutions preserve the rankings: documentary systems.}

Shift position and ask which poems survived before widespread printing. Those echoed by eminent writers, selected in major anthologies, and repeatedly copied into encyclopaedias and textbooks. Tao enjoyed two advantages: Liang crown prince Xiao Tong, compiler of Selections of Refined Literature, admired him, collected his works, and wrote a preface; after the Song, widespread printing fixed Su Shi's praise into countless editions. Another poet leaving little writing, lacking patrons and anthology inclusion, might lose even their name despite brilliance. This account restores luck and institutions, omitted by the others. Its limit: institutions preserve and disseminate but cannot manufacture love. If anthology inclusion guaranteed mastery, every writer in Selections would be a master; they are not. Institutions can determine who gets read, not who gets loved.

Each account grasps part of the truth. Remember the methodological warning: an explanation claiming everything has usually discarded something beforehand.

\section[Further Challenge: Beauty or Weapon?]{A Deeper Challenge: Is the Canon Beauty Itself, or a Weapon?}
Now introduce one of twentieth-century literary study's most famous disputes, representing two fundamentally different theories of the canon.

On one side stands American critic Harold Bloom. His influential 1994 The Western Canon discusses twenty-six supposedly unshakable canonical writers, from Shakespeare to Borges. His central claim is that \textbf{overwhelming aesthetic power makes a classic}: the work suddenly enlarges the world and makes language seem insufficient. He firmly opposed then-fashionable academic efforts to rearrange literary history politically, reallocating places by gender, race, and class. He acidly called these approaches the "School of Resentment", suggesting they settled scores with literature instead of reading it. You may dislike the canon, he argues, but cannot substitute votes for aesthetic judgement. Shakespeare became great not through power's promotion; power borrowed his radiance because he was great.

French sociologist Pierre Bourdieu can represent the other side. His 1979 Distinction advanced an uncomfortable claim: \textbf{supposedly refined taste is largely a class passport, not an inborn gift of appreciation}. Family concert visits, school-taught poems, a social circle's allusive jokes form cultural capital inherited like money, certified by schools and cashed in for prospects. The canon becomes a vast machine of cultural reproduction: the class possessing classics defines good taste, tests everyone's children against it, and declares its own children more cultured. Reading classics becomes learning identity passwords.

Weigh both carefully: they are genuine opponents, not simply right and wrong.

Bloom's strength is his stubborn defence of the work itself. If canonical status arose wholly from power, why do states so often fail to popularise their favoured authors? Every country has officially promoted "important works" promptly forgotten. Aesthetic experience is real: a bad book does not become enjoyable through textbook inclusion. His blind spot is treating aesthetics as if it occupied a vacuum, when upbringing, education, and class train our sense of beauty. Bourdieu asks precisely who fitted your aesthetic spectacles.

Bourdieu explains what Bloom cannot: why canonical lists heavily favour particular genders, classes, and languages, and why the ability to understand classics follows family background so neatly. His blind spot: if everything is cultural-capital competition, we cannot distinguish Dream of the Red Chamber from mediocre writing fortuitously promoted by power. Yet that distinction is real, as anyone who has genuinely read both knows.

Consider an \textbf{easily misread counterexample}. Czech writer Kafka worked for an insurance company, published a little, and found few sympathetic readers. Before dying in 1924, he instructed friend Max Brod to burn all unpublished manuscripts. Brod did not. He organised and published them, giving us The Trial, The Castle, and an indispensable twentieth-century name. Which sides does this challenge? "Works shine by themselves" fails: had Brod obeyed, Kafka's aesthetic power would be ashes, unable to shine however great. But "power manufactures every classic" fails too: Brod was no powerful authority, merely a reader convinced these manuscripts should live. No institution bestowed Kafka's status; generations of ordinary readers read it into being. The case shows that \textbf{canonisation is a relay: drop any baton and a great work cannot finish, but no runner can run the work's own leg for it}.

A similar twentieth-century relay had an identifiable next runner. African American writer Zora Neale Hurston wrote Their Eyes Were Watching God in the 1930s, presenting a Black woman's own voice in vernacular dialect. White mainstream literature neglected it, while some Black male critics complained it offered insufficient protest. Hurston died poor in a welfare home in 1960 and was buried without a headstone. Thirteen years later, young Black writer Alice Walker, later author of The Color Purple, travelled to find the unmarked grave and erected a stone. Through essays and edited collections she returned Hurston's books to readers one by one. Today Their Eyes Were Watching God is among the novels most often taught in American universities. Reviving a book sometimes literally requires a later reader to place a stone on a grave.

\section[Today: The New Bell Inside Your Phone]{Back to Today: The New Wooden-Tongued Bell Inside Your Phone}
This ancient question plays out in your pocket daily, with an entirely new cast.

Consider \textbf{online literature}. On platforms such as Qidian, visibility depends on clicks, monthly votes, tips, and algorithmic recommendations: no editor's manuscript judgement or critic's stamp, just readers voting with money and fingers. For the first time, gatekeepers are not intellectuals but a market--algorithm hybrid. Previously excluded people genuinely enter: a daytime delivery rider writing fantasy at night can acquire a million readers without a literature degree. Yet no checkpoint disappears; gatekeepers change. Luck becomes algorithmic discovery; publication becomes contracts and listing; criticism becomes ratings and chapter comments; teaching follows as online literature enters university courses and literary histories. The gate has become revolving, not vanished.

Consider \textbf{genre fiction's reversal of fortunes}. Martial-arts novels were once street-stall reading by despised literary hacks, yet Jin Yong's death brought the Chinese-speaking world's farewell to one of its era's most important writers. In 2004, a People's Education Press senior-secondary Chinese reader included Demi-Gods and Semi-Devils, provoking public debate over whether martial-arts fiction belonged in teaching. Science fiction, long dismissed as writing for children and engineers, received sudden serious mainstream attention after Liu Cixin's The Three-Body Problem won science fiction's highest honour, the Hugo, in 2015. Detective queen Agatha Christie ranks among the world's bestselling authors, yet literary history long allotted her only a corner. Gatekeepers build genre walls; readers on either side have never cared about them.

\textbf{Children's literature} reveals another double standard. Every child grows up with stories, yet literary history long confines their literature to a side room. It has its own gatekeeping institutions, such as America's Newbery Medal, awarded since 1922, but mainstream prizes rarely look its way. Harry Potter and the Philosopher's Stone, among history's bestselling children's novels, was rejected by more than a dozen publishers. Sometimes a childlike reader must correct gatekeepers' most candid failure.

\textbf{Literary prizes} also redraw literature's boundaries. Mo Yan's 2012 Nobel gave Chinese-language writing this stamp for the first time in a novelist's capacity. In 2016, American singer Bob Dylan won for new poetic expression within the song tradition, setting literary readers worldwide arguing: are lyrics literature? You already hold part of an answer. Homer was sung, the Book of Songs' Airs were songs, and West African griots still sing. Requiring literature to be prose printed on book pages is itself a young, local standard.

The fiercest disputes concern \textbf{textbooks}. Every Chinese syllabus revision makes news. Increases or decreases in Lu Xun's representation can trend for years. In Zhu Ziqing's "Retreating Figure", a father crosses railway platforms to buy oranges; online objections about breaking traffic rules sparked national laughter and debate. Seemingly trivial, these disputes make our question public. A syllabus alone is a list made by public authority and handed to every child. Everyone cares and has a right to address it. Argument brings gatekeeping out from backstage.

Beware a new species: \textbf{reading lists replacing reading}. Videos promising "One Hundred Years of Solitude in three minutes" or "Dream of the Red Chamber in one diagram" attract hundreds of millions of views. They shortcut the line, swallowing the evaluative authority of checkpoints two, three, and four, then announcing what a book means and whether you should read it. Italian writer Italo Calvino, who devoted an essay to classics, offered a famous formulation in Why Read the Classics?: classics are books people describe as rereading, not reading. It points to the fundamental distinction from bestsellers: a bestseller is consumed once; a classic is reopened across generations. Three-minute summaries pre-empt your first opening. After consuming someone else's conclusions, you forever lose your own first meeting with the book.

\section[Your Turn: Compile the Reading List]{Your Turn: What If You Compiled That Reading List?}
Return to the sheet outside the bookshop.

Suppose students a century hence must read ten books, and you alone must compile their list. All candidates stand on the shelves: thousand-year survivors, last year's sensations, languages you cannot read, unpublished manuscripts, online novels, songs, oral epics. By what criteria do you choose ten?

Before answering, put every weight already in your hands on the scales, enough on both sides.

If you say "the best-written", answer: what is good? Zhong Rong's ruler could not measure Tao Yuanming; can yours measure the next Tao? Bloom would say aesthetic experience, though hard to articulate, really exists, and people feigning confusion inwardly distinguish good books from bad. Dare you declare, as he did, that you possess a judgement and accept responsibility for it?

If you choose fairness, reserving places for different voices, remember there are only ten. Bourdieu warns that an uncorrected list merely carries old prejudice onward. Bloom asks how distributing places by representation differs methodologically from Qing-approved textbooks. Which is the greater loss: mediocre works selected fairly, or great works omitted through prejudice?

If you leave the choice to time and the market a century hence, remember Kafka's manuscripts stood one match from the fireplace and Hurston's grave lacked a stone for thirteen years. Time is not a judge. It consists of specific decisions by specific people, including you now, deciding whether to recommend this book to the next reader.

No standard answer exists. The list's dilemma miniaturises every humanistic choice: we cannot pretend to have no standards, which is dishonest, nor pretend our standards are impartial, which is naive. We can identify each list's gatekeepers, measuring sticks, and excluded people, then cast our own vote with greater awareness.

Next time you pass the bookshop, stay before that paper ten seconds longer. It is no longer paper but the mouth of a river flowing for two thousand years---and you stand in its sea.
